\section{Conclusions}\label{sec:conclusions}

We have presented \madsons, a new \mgshort\ module that automates the simulation of
arbitrary tree-level processes involving quarkonium and leptonium production.
Building upon the previous implementation of \swavetxt\ states~\cite{ColpaniSerri:2025vdz}, the present work completes the automation of the most phenomenologically relevant LO bound-state production processes by incorporating the orbital- and total-angular-momentum projectors required for \pwavetxt\ states.
The derivatives entering the orbital angular-momentum projector are evaluated using multivariate dual numbers, yielding exact and numerically stable results without relying on symbolic differentiation. 
The implementation has been extensively validated through comparisons with \helaconia. The squared helicity amplitudes agree for individual phase-space points at the level expected from double-precision arithmetic, while integrated cross sections are found to agree within the numerical MC uncertainties. 

We have further introduced a momentum-reshuffling procedure that reconciles the physical bound-state masses with the on-shell constituent masses required to evaluate the matrix elements. This treatment correctly reproduces the physical production thresholds and removes a limitation of conventional NRQCD calculations, particularly in processes involving several bound states composed of identical constituents but possessing different physical masses. Multiple reshuffling prescriptions, based on initial- and final-state recoil schemes, allow the residual ambiguity of the procedure to be quantified, providing a new source of theoretical uncertainty which is, in general, not negligible and should therefore be included in phenomenological studies.
The overall effect of including physical masses in the cross-section calculation is satisfactory. For most of the processes investigated in this work, our mass treatment reproduces the expected cross-section scaling, thereby improving the theoretical predictions. However, in some cases residual tensions remain, which may indicate that our implementation captures only part of the higher-order corrections in the $v$ expansion.

We illustrated the capabilities of our framework through single- and double-quarkonium production in proton-proton collisions at $\sqrt{s} = 13\,{\rm TeV}$ and charmonium production in $e^+ e^-$ collisions at $B$-factory energies. We also investigated the impact of using the measured quarkonium masses and quantified the associated reshuffling uncertainty. As a dedicated phenomenological application, we revisited $J/\psi + \psi(2\swave)$ production at the LHC, considering prompt $J/\psi$ production, including feed-down contributions from $\psi(2\swave)$ and $\chi_{cJ}$ states. The decay of the heavier states, parton showering, and hadronisation were simulated by interfacing the parton-level \mgshort\ events with \pythia{\tt 8}. Employing the physical quarkonium masses lowers the predicted cross section relative to the conventional treatment, leading to an improved description of the LHCb data while still leaving room for sizeable DPS contributions.

The framework also enables the study of leptonium production processes, providing a flexible tool for investigations of exotic atoms. In this context, we derived closed-form NRQED predictions for the considered cross sections and showed that \mgshort\ reproduces the analytic results for \pwavetxt\ para- and ortho-positronium production at the few-permille level, and for ortho-ditauonium production at the sub-permille level, across the full range of collision energies. This provides a fully independent validation of the \pwavetxt\ implementation. 

The \madsons\ module seamlessly integrates into the \mgshort\ framework, benefiting from its numerous existing features. Together with the previous \swavetxt\ implementation, this work enables the fully automated event generation of processes involving arbitrary numbers of quarkonia and leptonia within the established workflows. These developments provide a solid foundation for future extensions of \madsons, most notably towards the automation of NLO predictions within the same framework.